% part: intuitionistic-logic
% chapter: propositions-as-types
% section: proofs-to-terms
%READERNOTE{SOL6-A324: येथे दिलेल्या पद-नियमांसाठी सिद्धतेची विधानीय व्याप्ती स्पष्ट केला; संख्यापकांचे पद-नियम दिले असल्याचा दावा नाही.}

\documentclass[../../../include/open-logic-section]{subfiles}

\begin{document}

\olfileid{int}{pty}{pt}

\olsection{\usetoken{P}{derivation} चे सिद्धता-पदांत रूपांतर}

विधानीय नैसर्गिक निगमनातील !!{proof}s चे सिद्धता-पदांत रूपांतर,
म्हणजे~\Log{N2i} मधील !!{proof}s चे~\Log{tN2i}
मधील !!{proof}s मध्ये भाषांतर, अगदी सरळ आहे.
!!{derivation} मधील प्रत्येक क्रमवर्तीच्या उजव्या
बाजूच्या !!{formula} च्या डावीकडे सिद्धता-पद लिहायचे;
म्हणजे $\Gamma \Sequent M : !A$ या रूपातील
क्रमवर्ती मिळतात. मग संदर्भ~$\Gamma$ मध्ये~$M$ हे
$!A$ ची \emph{साक्ष देते} असे म्हणू. कोणते सिद्धता-पद
लिहायचे हे पूर्णपणे \Log{tN2i} च्या नियमांनी ठरते.

\begin{prop}
विधानीय \Log{N2i} मधील $\Gamma \Sequent !A$ ची प्रत्येक
!!{proof}~$\delta$ ही~\Log{tN2i} मधील
$\Gamma \Sequent N : !A$ च्या !!a{proof}~$\delta'$ शी
अनुरूप असते. सिद्धता-पद~$N$ हे~$\delta$ वरून
एकमेवपणे ठरते.
\end{prop}

\begin{proof}
~$\delta$ च्या उंचीवर विगमन करू. $\delta$ हे
$x:!A \Sequent !A$ हे स्वयंसिद्धक असेल, तर
$\delta'$ हे~$x:!A \Sequent x:!A$ हे स्वयंसिद्धक आहे.
$\delta$ एखाद्या अनुमान-नियमावर संपत असेल, तर
विगमन-गृहीतक प्रत्येक आधारविधानासाठी सिद्धता-पद आणि
ते पद असलेल्या क्रमवर्तीची \Log{tN2i}-!!{proof}s देते.
मग \Log{tN2i} चा अनुरूप नियम लावतो;
\Log{tN2i} चा अनुमान-नियमच अनुरूप सिद्धता-पद निश्चित करतो.
\end{proof}
येथे दिलेल्या पद-नियमांची व्याप्ती विधानीय आहे;
संख्यापकांचे स्वतंत्र पद-संकारक परिभाषित केलेले नाहीत.

\begin{ex}
  \Log{N1i} मधील $(!A \land !B) \lif (!B \land !A)$
  ची पुढील अगदी साधी !!{proof} घ्या:
  \[
  \AxiomC{$\Discharge{!A\land !B}{x}$}
  \RightLabel{\Elim\land[2]}
  \UnaryInfC{$!B$}
  \AxiomC{$\Discharge{!A\land !B}{x}$}
  \RightLabel{\Elim\land[1]}
  \UnaryInfC{$!A$}
  \RightLabel{\Intro\land}
  \BinaryInfC{$!B \land !A$}
  \DischargeRule{\Intro\lif}{x}
  \UnaryInfC{$(!A \land !B) \lif (!B \land !A)$}
  \DisplayProof
  \]
  \Log{N2i} मध्ये ती अशी दिसते:
  \[
  \Axiom$x : !A\land !B \fCenter !A\land !B$
  \RightLabel{\Elim\land[2]}
  \UnaryInf$x : !A\land !B \fCenter !B$
  \Axiom$x : !A\land !B \fCenter !A\land !B$
  \RightLabel{\Elim\land[1]}
  \UnaryInf$x : !A\land !B \fCenter !A$
  \RightLabel{\Intro\land}
  \BinaryInf$x : !A\land !B \fCenter !B \land !A$
  \RightLabel{\Intro\lif}
  \UnaryInf$\fCenter (!A \land !B) \lif (!B \land !A)$
  \DisplayProof
  \]
  सिद्धता-पदे वरपासून खाली देऊ. स्वयंसिद्धकांची
  सिद्धता-पदे ही संदर्भातील~$x$ या चलासारखीच आहेत.
  मग $\Elim\land[i]$ शी संबंधित सिद्धता-पदे अनुक्रमे
  $\proj{2}{x}$ आणि~$\proj{1}{x}$ अशी ठरतात.
  \Intro{\land} लावताना ही दोन सिद्धता-पदे एकत्र करतो:
  $\pair{\proj{2}{x}}{\proj{1}{x}}$.
  शेवटी \Intro{\lif} नियमाने $\lambd$-अमूर्तीकरण
  केले जाते; त्यामुळे~$x$ बद्ध होतो आणि~$x$ ही खूण
  $!A \land !B$ या गृहीतकाला होती हे नोंदले जाते:
  $\lambd[\typeof{x}{!A \land
  !B}][\pair{\proj{2}{x}}{\proj{1}{x}}]$.
  मग \Log{tN2i} मधील !!{proof} अशी आहे:
  \[
  \Axiom$x : !A\land !B \fCenter x: !A\land !B$
  \RightLabel{\Elim\land[2]}
  \UnaryInf$x : !A\land !B \fCenter \proj{2}{x} : !B$
  \Axiom$x : !A\land !B \fCenter x: !A\land !B$
  \RightLabel{\Elim\land[1]}
  \UnaryInf$x : !A\land !B \fCenter \proj{1}{x} : !A$
  \RightLabel{\Intro\land}
  \BinaryInf$x : !A\land !B \fCenter \pair{\proj{2}{x}}{\proj{1}{x}} :
    !B \land !A$
  \RightLabel{\Intro\lif}
  \UnaryInf$\fCenter \lambd[\typeof{x}{!A \land
  !B}][\pair{\proj{2}{x}}{\proj{1}{x}}] : (!A \land !B) \lif (!B \land !A)$
  \DisplayProof
  \]
\end{ex}

\end{document}
